\documentclass[11pt, aspectratio=169]{beamer}
\usetheme{Modern}
\usepackage{pgfplots}
\usepackage{amssymb}
\usepackage{booktabs}
\usepackage{tikz}
\usetikzlibrary{trees,shapes,arrows,positioning,decorations.pathreplacing}
\usepackage[backend=biber,style=numeric,sorting=none]{biblatex}
\addbibresource{literature.bib}
\pgfplotsset{compat=1.18}

\title{Maschinelles Lernen}
\subtitle{Eine kurze Einführung für JuristInnen}
\author{Julian Streitberger (FMA -- Abteilung IV/4)}
\date{\today}

\begin{document}

\maketitle

%===============================================================================
% FOLIE 1: Was ist ML?
%===============================================================================

\begin{frame}
    \frametitle{Was ist Maschinelles Lernen?}
    \begin{columns}
        \begin{column}{0.55\textwidth}
            \begin{block}{Kernidee}
                Statt Regeln manuell zu programmieren, lernt der Computer die Regeln \textbf{aus Daten}.
            \end{block}
            \vspace{0.4cm}
            \textbf{Traditionelle Programmierung:}\\
            Mensch schreibt Regeln $\rightarrow$ Computer führt aus
            \vspace{0.3cm}
            
            \textbf{Maschinelles Lernen:}\\
            Daten + Beispiele $\rightarrow$ Computer lernt Regeln
        \end{column}
        \begin{column}{0.42\textwidth}
            \begin{exampleblock}{Beispiel: Geldwäsche}
                \textbf{Regelbasiert:}\\
                „Melde alle TX $>$ \texteuro 10.000"
                \vspace{0.2cm}
                
                \textbf{ML-basiert:}\\
                System erkennt verdächtige Muster aus historischen Fällen -- auch \textit{unbekannte}.
            \end{exampleblock}
        \end{column}
    \end{columns}
\end{frame}

%===============================================================================
% FOLIE 2: Zwei Lernarten
%===============================================================================

\begin{frame}
    \frametitle{Zwei Lernarten: Überwacht vs. Unüberwacht}
    \begin{columns}
        \begin{column}{0.48\textwidth}
            \begin{block}{Überwachtes Lernen (Supervised)}
                \begin{itemize}
                    \item Computer lernt aus \textbf{Beispielen mit Lösung}
                    \item Wie ein Student mit Musterlösungen
                    \item \textit{Bsp: „Diese 1000 Fälle waren Betrug, diese 9000 nicht"}
                \end{itemize}
                \vspace{0.2cm}
                \textbf{Anwendung:} Alert-Klassifikation, Risiko-Scoring
            \end{block}
        \end{column}
        \begin{column}{0.48\textwidth}
            \begin{block}{Unüberwachtes Lernen (Unsupervised)}
                \begin{itemize}
                    \item Computer findet \textbf{selbst Muster}
                    \item Wie ein Detektiv ohne Hinweise
                    \item \textit{Bsp: „Gruppiere ähnliche Kunden"}
                \end{itemize}
                \vspace{0.2cm}
                \textbf{Anwendung:} Kundensegmentierung, Anomalieerkennung
            \end{block}
        \end{column}
    \end{columns}
\end{frame}

%===============================================================================
% FOLIE 3: Das Problem Overfitting
%===============================================================================

\begin{frame}
    \frametitle{Das zentrale Problem: Overfitting}
    \begin{center}
        \includegraphics[width=0.85\linewidth]{plots/overfitting_visualization.pdf}
    \end{center}
    \begin{alertblock}{Juristische Analogie}
        \textbf{Overfitting} = Das Modell „lernt den Fall auswendig" statt das Prinzip zu verstehen.\\
        Wie ein Jurist, der nur die Urteile kennt, aber nicht die dahinterliegenden Rechtsgrundsätze.
    \end{alertblock}
\end{frame}

%===============================================================================
% FOLIE 4: Metriken
%===============================================================================

\begin{frame}
    \frametitle{Wie gut ist ein Modell? -- Precision \& Recall}
    \begin{columns}
        \begin{column}{0.55\textwidth}
            \textbf{Das Dilemma bei seltenen Ereignissen:}
            \begin{itemize}
                \item Geldwäsche: $<$1\% aller Transaktionen
                \item Ein Modell, das \textit{immer} „unbedenklich" sagt, hat 99\% Trefferquote!
            \end{itemize}
            \vspace{0.3cm}
            \textbf{Bessere Maßstäbe:}
            \begin{itemize}
                \item \textbf{Precision:} Wie viele Alarme sind echte Treffer?
                \item \textbf{Recall:} Wie viele echte Fälle wurden gefunden?
            \end{itemize}
        \end{column}
        \begin{column}{0.42\textwidth}
            \begin{alertblock}{Abwägung}
                \textbf{Hoher Recall:}\\
                Weniger übersehene Fälle\\
                {\small (regulatorisch wichtig!)}
                \vspace{0.2cm}
                
                \textbf{Hohe Precision:}\\
                Weniger Fehlalarme\\
                {\small (Effizienz)}
            \end{alertblock}
        \end{column}
    \end{columns}
\end{frame}

%===============================================================================
% FOLIE 5: Fragen für Prüfer
%===============================================================================

\begin{frame}
    \frametitle{Fragen für Vor-Ort-PrüferInnen}
    \begin{columns}
        \begin{column}{0.48\textwidth}
            \begin{block}{Zur Modellentwicklung}
                \begin{itemize}
                    \item Mit welchen \textbf{Daten} wurde trainiert?
                    \item Wurde auf \textbf{Overfitting} getestet?
                    \item Welche \textbf{Faktoren} fließen ein?
                \end{itemize}
            \end{block}
            \vspace{0.2cm}
            \begin{block}{Zur Validierung}
                \begin{itemize}
                    \item Welche \textbf{Metriken} werden gemessen?
                    \item Wie oft wird \textbf{überprüft}?
                \end{itemize}
            \end{block}
        \end{column}
        \begin{column}{0.48\textwidth}
            \begin{block}{Zur Governance}
                \begin{itemize}
                    \item Wer ist \textbf{verantwortlich}?
                    \item Ist das Modell \textbf{erklärbar}?
                    \item Wie werden \textbf{Änderungen} dokumentiert?
                \end{itemize}
            \end{block}
            \vspace{0.2cm}
            \begin{alertblock}{Red Flags}
                \begin{itemize}
                    \item Keine Dokumentation
                    \item Nur „Accuracy" als Metrik
                    \item Kein regelmäßiges Re-Training
                \end{itemize}
            \end{alertblock}
        \end{column}
    \end{columns}
\end{frame}

%===============================================================================
% FOLIE 6: Key Takeaways
%===============================================================================

\begin{frame}
    \frametitle{Zusammenfassung}
    \begin{enumerate}
        \item \textbf{ML lernt aus Daten} -- keine explizite Programmierung nötig
        \vspace{0.3cm}
        \item \textbf{Überwacht vs. Unüberwacht:} Mit oder ohne Beispiellösungen
        \vspace{0.3cm}
        \item \textbf{Overfitting vermeiden:} Das Modell soll generalisieren können
        \vspace{0.3cm}
        \item \textbf{Precision \& Recall:} Wichtiger als reine Trefferquote
        \vspace{0.3cm}
        \item \textbf{Governance \& Erklärbarkeit:} Essenziell für regulierte Bereiche
    \end{enumerate}
    \vspace{0.4cm}
    \begin{block}{Literaturempfehlung}
        Géron, A. (2022). \textit{Hands-On Machine Learning}. O'Reilly. \cite{geron2022hands}
    \end{block}
\end{frame}

\begin{frame}[allowframebreaks]
    \frametitle{Literatur}
    \printbibliography[heading=none]
\end{frame}

\end{document}